\documentclass[11pt]{article}
\usepackage[a4paper,margin=27mm]{geometry}
\usepackage{amsmath,amssymb,amsthm,hyperref}
\newtheorem{lemma}{Lemma}
\newtheorem{theorem}{Theorem}
\title{Six weak triple-minimum permutations support at most six labels}
\author{Research note: a proof independent of constraint solvers}
\date{30 September 2026}
\begin{document}
\maketitle
All permutation families below are multisets of six equally weighted
permutations. A weak triple-minimum family has each member of each
triple earliest in exactly two rows. A fully triple-uniform family
has every ordering of every triple in exactly one row.

\begin{lemma}\label{psca}
No fully triple-uniform six-row family exists on five labels.
\end{lemma}
\begin{proof}
For a fixed label let $d_j$ be its number of occurrences in position
$j\in\{0,1,2,3,4\}$. Counting labels before it, and pairs of labels
before it, gives
\[
 \sum_jd_j=6,\qquad \sum_jj d_j=12,\qquad
 \sum_jj(j-1)d_j=24.
\]
The last identity modulo three implies $3\mid d_2$.
The possibility $d_2=6$ contradicts $\sum j^2d_j=36$.
If $d_2=3$, put $(a,b,c,d)=(d_0,d_1,d_3,d_4)$. The remaining
equations are
\[
 a+b+c+d=3,\qquad b+3c+4d=6,\qquad b+9c+16d=24.
\]
Thus $c+2d=3$ and $b=2d-3$. Since $c\ge0$ forces $d\le1$,
this makes $b<0$. Therefore $d_2=0$ for every label, impossible
because every row has a label in position two.
\end{proof}
Lemma~\ref{psca} is the known nonexistence of
$\mathrm{PSCA}(5,3,1)$; the position-moment proof is also given in
Gentle and Wanless \cite{GW}.

\begin{theorem}
A weak triple-minimum family of six permutations has at most six
labels. The bound is attained.
\end{theorem}
\begin{proof}
It suffices to exclude seven labels, since deleting labels preserves
the property. For distinct labels $i,j$ let $r_{ij}$ count rows in
which $i$ precedes $j$. Triple-minimum uniformity implies
$2\le r_{ij}\le4$.

Let $B$ be the labels $i$ for which $r_{ij}=2$ for some $j$, and let
$A$ be the other labels. If $i\in B$, the two rows in which $i<j$
must also have $i<k$ for every third label $k$, because the latter
joint event already occurs twice. Thus $i$ is globally first in
exactly two rows. There can be at most three such labels, so
$|A|\ge4$.

For any two labels of $A$, neither order count can be two, hence both
are three. On a triple, equal minimum counts together with balanced
pair counts force all six orders to occur once. Consequently the
restriction to $A$ is fully triple-uniform. Lemma~\ref{psca} gives
$|A|\le4$, so $|A|=4$ and $|B|=3$. Each member of $B$ is first in
two rows, accounting for all six first positions.

Fix $a\in A$. Sum its minimum counts over the three triples
$\{a,b,c\}$ with $b,c\in B$. The sum is six. In any row at least
one member of $B$ precedes $a$, so that row contributes at most one
to this sum. It follows that every row contributes one: exactly one
member of $B$ precedes $a$. Thus in each row all four labels in $A$
occur between the first member of $B$ and the other two members.

Fix $b\in B$ and distinct $a,c\in A$. Outside the two rows beginning
with $b$, both $a,c$ precede $b$. The minimum condition on
$\{a,c,b\}$ therefore says that $a<c$ holds in exactly two of those
four rows. Since it holds in three rows overall, it holds in exactly
one of the two rows beginning with $b$. This holds for every pair
from $A$, so the restrictions of those two rows to $A$ are reverses.
The fully triple-uniform family on $A$ is therefore three reversal
pairs.

Choose one representative of each reversal pair. For every triple of
$A$, every label must be its middle label in exactly one of these
three representatives. A fixed label of $A$ would consequently be
middle in three triples in total. But in any permutation on four
labels, a label at position $j\in\{0,1,2,3\}$ is middle in exactly
$j(3-j)\in\{0,2\}$ triples. The sum over three representatives is
even, a contradiction.

Sharpness is witnessed by the six permutations
\[
 012345,\quad521340,\quad531204,\quad413250,\quad423105,\quad032145,
\]
whose triple-minimum counts are checked in
\path{verify_nine_dice_cloning.py}.
\end{proof}
This proof concerns six equally weighted rows. The stronger weighted
obstruction is now established in
\path{weak_minimum_equality_obstruction.tex}. Neither result excludes
ten twelve-face three-wise fair dice by another construction.
No claim of literature novelty is made here.

\begin{thebibliography}{1}
\bibitem{GW} B. Gentle and I. M. Wanless,
\emph{On Perfect Sequence Covering Arrays},
Annals of Combinatorics (2023).
\href{https://arxiv.org/abs/2202.01960}{arXiv:2202.01960},
\href{https://doi.org/10.1007/s00026-022-00610-6}{doi:10.1007/s00026-022-00610-6}.
\end{thebibliography}
\end{document}
